% ENEM 2016-2022 — 5 questões MCQ — Função Quadrática
% Fonte: lista "Função do 2º grau ou Função Quadrática" do site projetoagathaedu.com.br
% Omitidas: Q1 (imagem), Q2 (duplicata lista14/Q1), Q8,Q9,Q10,Q11,Q12 (imagens)
%           Q6 (figura define o contexto mas resposta é conceitual — incluída)

---
tipo: Múltipla Escolha
dificuldade: Fácil
nivel: médio
disciplina: Matemática
assunto: Função Quadrática
gabarito: C
tags: [epidemia, infectados, vértice, propaganda, mês]
fonte: concurso
concurso: ENEM
banca: INEP
ano: 2022
numero: 3
---
\question
Ao analisar os dados de uma epidemia em uma cidade, peritos obtiveram um modelo que avalia a quantidade de pessoas infectadas a cada mês, ao longo de um ano. O modelo é dado por \( p(t) = -t^2 + 10t + 24 \), sendo \( t \) um número natural, variando de 1 a 12, que representa os meses do ano, e \( p(t) \) a quantidade de pessoas infectadas no mês \( t \) do ano. Para tentar diminuir o número de infectados no próximo ano, a Secretaria Municipal de Saúde decidiu intensificar a propaganda oficial sobre os cuidados com a epidemia. Foram apresentadas cinco propostas (I, II, III, IV e V), com diferentes períodos de intensificação das propagandas:

-- I: \( 1 \leq t \leq 2 \);

-- II: \( 3 \leq t \leq 4 \);

-- III: \( 5 \leq t \leq 6 \);

-- IV: \( 7 \leq t \leq 9 \);

-- V: \( 10 \leq t \leq 12 \).

A sugestão dos peritos é que seja escolhida a proposta cujo período de intensificação da propaganda englobe o mês em que, segundo o modelo, há a maior quantidade de infectados. A sugestão foi aceita.

A proposta escolhida foi a
\begin{oneparchoices}
\choice I.
\choice II.
\correctchoice III.
\choice IV.
\choice V.
\end{oneparchoices}

---
tipo: Múltipla Escolha
dificuldade: Fácil
nivel: médio
disciplina: Matemática
assunto: Função Quadrática
gabarito: A
tags: [fluido, tubo, velocidade, escoamento, máximo]
fonte: concurso
concurso: ENEM
banca: INEP
ano: 2021
numero: 4
---
\question
Considere que o modelo matemático utilizado no estudo da velocidade \( V \), de uma partícula de um fluido escoando em um tubo, seja diretamente proporcional à diferença dos quadrados do raio \( R \) da secção transversal do tubo e da distância \( x \) da partícula ao centro da secção que a contém. Isto é, \( V(x) = K^2(R^2 - x^2) \), em que \( K \) é uma constante positiva.

O valor de \( x \), em função de \( R \), para que a velocidade de escoamento de uma partícula seja máxima é de
\begin{oneparchoices}
\correctchoice 0.
\choice \( R \)
\choice \( 2R \).
\choice \( KR \).
\choice \( K^2 R^2 \).
\end{oneparchoices}

---
tipo: Múltipla Escolha
dificuldade: Difícil
nivel: médio
disciplina: Matemática
assunto: Função Quadrática
gabarito: E
tags: [navios, codificação, operações, equação quadrática]
fonte: concurso
concurso: ENEM
banca: INEP
ano: 2021
numero: 5
---
\question
Para a comunicação entre dois navios é utilizado um sistema de codificação com base em valores numéricos. Para isso, são consideradas as operações triângulo \( \Delta \) e estrela \( * \), definidas sobre o conjunto dos números reais por \( x \Delta y = x^2 + xy - y^2 \) e \( x * y = xy + x \).

O navio que deseja enviar uma mensagem deve fornecer um valor de entrada \( b \), que irá gerar um valor de saída, a ser enviado ao navio receptor, dado pela soma das duas maiores soluções da equação \( (a \Delta b) * (b \Delta a) = 0 \). Cada valor possível de entrada e saída representa uma mensagem diferente já conhecida pelos dois navios.

Um navio deseja enviar ao outro a mensagem "ATENÇÃO!". Para isso, deve utilizar o valor de entrada \( b = 1 \).

Dessa forma, o valor recebido pelo navio receptor será
\begin{oneparchoices}
\choice \( \sqrt{5} \)
\choice \( \sqrt{3} \)
\choice \( \sqrt{1} \)
\choice \( \dfrac{-1+\sqrt{5}}{2} \)
\correctchoice \( \dfrac{3+\sqrt{5}}{2} \)
\end{oneparchoices}

---
tipo: Múltipla Escolha
dificuldade: Média
nivel: médio
disciplina: Matemática
assunto: Função Quadrática
gabarito: B
tags: [andaime, trajetória, domínio, raízes, queda]
fonte: concurso
concurso: ENEM-PPL
banca: INEP
ano: 2022
numero: 6
---
\question
A trajetória de uma pessoa que pula de um andaime até o chão é descrita por uma função \( y = f(x) \), sendo \( x \) e \( y \) medidos em metro. Seja \( D \) o domínio da função \( f(x) \).

Para que a situação representada seja real, o domínio dessa função deve ser igual a
\begin{choices}
\choice \( \{x_2\} \), sendo \( x_2 \) a raiz positiva de \( f(x) \).
\correctchoice \( \{x \in \mathbb{R} \mid 0 \leq x \leq x_2\} \), sendo \( x_2 \) a raiz positiva de \( f(x) \).
\choice \( \{x \in \mathbb{R} \mid x_1 \leq x \leq x_2\} \), sendo \( x_1 \) e \( x_2 \) raízes de \( f(x) \), com \( x_1 < x_2 \).
\choice \( \{x \in \mathbb{R} \mid x \geq 0\} \)
\choice \( x \in \mathbb{R} \).
\end{choices}

---
tipo: Múltipla Escolha
dificuldade: Média
nivel: médio
disciplina: Matemática
assunto: Função Quadrática
gabarito: B
tags: [dengue, epidemia, dedetização, infectados, equação]
fonte: concurso
concurso: ENEM
banca: INEP
ano: 2016
numero: 7
---
\question
Para evitar uma epidemia, a Secretaria de Saúde de uma cidade dedetizou todos os bairros, de modo a evitar a proliferação do mosquito da dengue. Sabe-se que o número \( f \) de infectados é dado pela função \( f(t) = -2t^2 + 120t \) (em que \( t \) é expresso em dia e \( t = 0 \) é o dia anterior à primeira infecção) e que tal expressão é válida para os 60 primeiros dias da epidemia.

A Secretaria de Saúde decidiu que uma segunda dedetização deveria ser feita no dia em que o número de infectados chegasse à marca de 1.600 pessoas, e uma segunda dedetização precisou acontecer.

A segunda dedetização começou no
\begin{oneparchoices}
\choice 19º dia.
\correctchoice 20º dia.
\choice 29º dia.
\choice 30º dia.
\choice 60º dia.
\end{oneparchoices}
